\chapter{Ecological, Cultural, and Computational Projections}
\label{ch:ecological_cultural}

\epigraph{Across ecological, cultural, and artificial domains, the invariant dynamic
is stigmergic reinforcement coupled to Markov blanket repair.}{}

\section{The Spectrum from Squirrel to Language Model}

The Yarncrawler principle is that agents do not merely consume; they weave, cache, and
repair. Each action leaves a trace that both satisfies immediate needs and sets the
stage for recursive self-maintenance. This dynamic recurs across ecological, cultural,
and computational scales. The invariants under this scale change are the \RSVP field
correspondences: $\Phi$ measures accumulated deposits, $\mathbf{v}$ encodes recursive
flows, and $S$ quantifies robustness.

\section{Ecological Yarncrawlers}

A squirrel does not carry a metric map of its forest. Its survival depends on weighted
affordances: gradients of safety, concealment, and food density blurred across its
perceptual field. When hungry, it follows Gaussian-blurred gradients toward likely
caches rather than performing global optimization. Its over-provisioning of seeds is
not an error but the creation of semantic attractors---distributed caches that stabilize
its future viability.

The perceived affordance field is $\psi(x) = (G_\sigma * u)(x)$ where $u(x)$ is the
distribution of hidden resources and $G_\sigma$ is a Gaussian smoothing kernel.
Foraging dynamics are $\dot{x}(t) = -\nabla\psi(x) + \eta(t)$. Each cache deposit
increments $u(x_k)$, making the field the squirrel uses simultaneously a navigation
tool and a stigmergic memory. The forest itself becomes a distributed memory field,
with excess ensuring robustness against local failure.

\section{Cultural Yarncrawlers}

Cultural Yarncrawlers instantiate the same principles at a larger scale. In Amazonian
communities, household waste was exported to boundary piles that grew fertile through
microbial and chemical processes. Once above a critical threshold, these berms
self-maintained, attracting further deposition and forming Voronoi-like tessellations
of territory. The resulting terra preta soils persist for millennia.

The governing dynamics are analyzed in the companion paper \emph{Yarncrawler in
Action} through a system of reaction-advection-diffusion equations for berm biomass
$B$, char matrix $C$, and nutrients $N$. The stigmergic threshold condition
$\alpha_B \langle S \rangle_{A_\tau} > \delta_B \theta_{\rm eff}$ formalizes the
transition from fragile deposition to self-maintaining growth.

\section{Artificial Yarncrawlers}

Mixture-of-Experts architectures and retrieval-augmented generation instantiate the
Yarncrawler principle computationally. Each expert module functions as a local chart
on the semantic manifold. Retrieval mechanisms act as stigmergic glue: a memory trace
that is frequently retrieved becomes more salient, just as a cairn or berm attracts
additional reinforcement. The global semantic field is assembled from local expert
densities by $\Phi(x) = \sum_i w_i(x)\varphi_i(x)$.

The key distinction from standard MoE/RAG architectures is that the repair mechanisms
of a Yarncrawler reconfigure the modules themselves rather than only gating access to
them. The semantic structure is continuous and trajectory-based, modeled by sheaf
gluing and \RSVP fields, rather than discrete or static.

\section{The Climate Response as Worked Example}

The read-evaluate-write-move loop provides the clearest operational instantiation of
the \KES map $\Omega_t \to H_{t+1}$ anywhere in the corpus. Consider a Yarncrawler
deployed across a coastal region:

\medskip
\noindent\textbf{Node 1 (Seattle):} Reads rising sea levels and failing levees.
Evaluates that current barriers are underfunded. Writes a plan to redirect local
taxes to infrastructure. Moves to Node 2, carrying the tax-redirection strategy
as a candidate fix for other urban nodes.

\noindent\textbf{Node 2 (Olympic National Forest):} Reads increased wildfire risk.
Evaluates that invasive species are worsening fire behavior. Writes a schedule for
controlled burns and native replanting. Moves to Node 3.

\noindent\textbf{Node 3 (Portland):} Reads heatwave data. Evaluates that cooling
centers are overcrowded. Writes a policy to expand public cooling spaces, using
lessons from Seattle's resource allocation. Moves to Node 4.
\medskip

The Yarncrawler learns that tax redirection from Seattle outperforms Portland's
grant-based model and applies this to future urban nodes. It notes that controlled
burns require longer monitoring cycles, refining its repair strategy for other
forest nodes. The trail of fixes constitutes a semantic ledger of the reasoning
trajectory---what the companion framework calls Chain of Memory.

\section{Forests as Distributed Constraint Networks}
\label{sec:forests}

\claimP Forests may be understood not as collections of isolated organisms but as
dynamically coupled systems in which constraint propagation occurs across multiple
interacting substrates simultaneously. Nutrient flow, signaling gradients, fungal
mediation, root competition, hydrological coupling, and atmospheric exchange
together generate a distributed regulatory field whose structure evolves over time
through repeated interaction.

Large persistent trees function not as centrally intelligent agents issuing
commands but as highly connected constraint hubs embedded within long-lived
ecological trajectories. Their historical persistence grants them unusually large
influence over local energetic accessibility conditions, nutrient distributions,
and signaling pathways. The resulting forest architecture behaves less like a
hierarchy of independent individuals and more like a partially shared admissibility
manifold in which the actions of one region alter the reachable futures of
neighboring regions.

\subsection{Kin Recognition as Constraint Compatibility}

Observed forms of kin recognition in forests need not imply anthropomorphic
intentionality. They may instead be interpreted as local compatibility detection
across coupled biochemical and structural histories. Root exudates act as
persistent chemical signatures encoding developmental and metabolic history.

\claimP Let $K(x_i,x_j)$ represent compatibility between organisms $x_i$ and
$x_j$. Rather than functioning as symbolic recognition, compatibility emerges from
overlap in biochemical trajectory structure:
\[
  K(x_i,x_j) = \exp\bigl(-\alpha D(H_i,H_j)\bigr),
\]
where $H_i$, $H_j$ denote historical biochemical configurations and $D$ measures
divergence between developmental trajectories. High compatibility corresponds to
reduced energetic conflict and increased exchange stability. Kin recognition
becomes a consequence of reduced incompatibility across historically related
developmental systems rather than explicit semantic identification.

\subsection{Mycorrhizal Networks as Ecological Routing Structures}

\claimP Mycorrhizal fungal systems may be interpreted as distributed transport
and signaling infrastructures that dynamically route energetic and informational
gradients throughout ecological space. The fungal substrate forms a continuously
adapting conductive geometry whose topology evolves in response to nutrient
pressure, environmental instability, and local metabolic demand.

Let the ecological transport graph be $G = (V,E)$, where vertices represent
organisms and edges represent fungal transport pathways. Resource transfer between
nodes can be modeled as constrained diffusion:
\[
  \frac{dR_i}{dt} = \sum_j w_{ij}(R_j - R_i) - \lambda_i R_i,
\]
where $R_i$ denotes local resource density and $w_{ij}$ encodes effective
conductive coupling strength. The network redistributes energetic accessibility
according to dynamic structural conditions rather than fixed optimization targets.

\subsection{Preferential Allocation and Developmental Stabilization}

\claimP Preferential transfer of carbon and nutrients toward younger saplings may
be interpreted as stabilization of nearby developmental trajectories. Young
organisms occupy fragile regions of ecological state space characterized by high
instability and low energetic accessibility. Transfers from persistent mature trees
locally reduce developmental volatility by expanding the reachable future states
available to saplings during early growth.

Define developmental stability as $S_d(x,t) = 1/(1+\sigma(x,t))$ where
$\sigma(x,t)$ measures local environmental volatility. Resource transfer functions
then preferentially amplify trajectories with low stability:
\[
  T_{ij} = K(x_i,x_j) \cdot \frac{1}{S_d(x_j,t)+\epsilon}.
\]
Persistent ecological structures naturally redistribute resources toward unstable
neighboring trajectories whose collapse would reduce overall network continuity.

\subsection{Ecological Memory and Structural Residue}

\claimP Forests accumulate persistent structural residue across timescales far
beyond individual organisms. Nutrient pathways, fungal geometries, soil chemistry,
canopy distributions, and hydrological routing collectively encode traces of prior
ecological interaction. This is ecological memory physically embedded within
evolving substrate---not symbolic or representational but infrastructural.
Previous interactions reshape the future accessibility structure of the environment
itself. Ecological intelligence, on this reading, emerges not from centralized
cognition but from long-horizon stabilization across recursively interacting
constraint systems.

\begin{constraintboundary}
\textbf{Status:~\textbf{[P]}} These interpretations are phenomenological
applications of the trajectory-admissibility framework to ecological systems.
They are supported by empirical forest ecology and mycorrhizal network research
but are not derived from those fields.

\textbf{Not claimed:} That forests possess consciousness, that kin recognition
is semantic, or that mycorrhizal networks function as computational systems in
the technical sense.

\textbf{Falsification:} A demonstration that nutrient redistribution in forest
networks is better explained by simple diffusion gradients without compatibility
modulation would weaken the constraint-compatibility interpretation of kin
recognition. Absence of preferential transfer toward low-stability saplings
under controlled conditions would weaken the developmental stabilization claim.
\end{constraintboundary}
